\section{Template C++17}

Base com\'un para los snippets del handbook. Incluye el template cl\'asico del
equipo; las macros son opcionales, pero permiten pegar c\'odigo viejo sin
traducir nombres durante una competencia.

\begin{lstlisting} [style=cpp]
#include <bits/stdc++.h>
#define pb push_back
#define fst first
#define snd second
#define fore(i,a,b) for (int i = (a), ThxDem = (b); i < ThxDem; ++i)
#define SZ(x) ((int)(x).size())
#define ALL(x) (x).begin(), (x).end()
#define mset(a,v) memset((a), (v), sizeof(a))
#define FIN ios::sync_with_stdio(false); cin.tie(nullptr); cout.tie(nullptr)
using namespace std;

using ll = long long;
using ull = unsigned ll;
using i128 = __int128_t;
using u128 = __uint128_t;
using ii = pair<int,int>;
using pii = pair<int,int>;
using pll = pair<ll,ll>;

const int INF = numeric_limits<int>::max() / 4;
const ll LINF = numeric_limits<ll>::max() / 4;

int main() {FIN;

    return 0;
}
\end{lstlisting}

Para llenar memoria byte a byte, usar \texttt{memset} solamente con valores
seguros como cero o $-1$. Para otros valores, usar \texttt{fill}.

\warning{\texttt{fore} guarda el l\'imite una sola vez: no usarlo si ese l\'imite
debe cambiar dentro del loop. (Es decir, no se recalcula con cada iteraci\'on como si lo har\'ia i < SZ(v) \texttt{pb/fst/snd/ALL/SZ} son abreviaturas; ante
un error de tipos, expandirlas mentalmente antes de depurar.}
